\chapter{Overview}\label{s:overview}

This chapter presents the system built to address the problem stated in
Chapter~\ref{s:intro}: translating a natural-language scientific intent into a
correct, executable BraneScript workflow. We describe the system from the
perspective of a user submitting an intent, then walk through the pipeline
stages that turn that intent into a validated, executed workflow.

\section{User Perspective}

A researcher interacts with the system by submitting a single natural-language
sentence, for example: \emph{``I want to analyse all heart-disease data using
package Healthcare''}. No knowledge of BraneScript syntax, package function
signatures, or Brane's execution model is required. The system returns one of
three outcomes: (1)~a BraneScript workflow together with its execution result,
(2)~a clarification question if the intent is ambiguous or underspecified
(e.g., no target dataset is named), or (3)~a final failure report after
repeated generation attempts could not produce a working workflow.

This interaction is available through two entry points that share the same
underlying pipeline: a command-line invocation of \texttt{pipeline.py}, and the
``Generate'' tab of the project's web dashboard, which submits the same request
as a batch job on the Snellius HPC cluster and streams back progress and results
(Section~\ref{s:implementation}).

\section{System Architecture}

Figure~\ref{fig:pipeline-architecture} shows the end-to-end architecture of the
system, from intent to executed workflow.

\begin{figure}[htbp]
    \centering
    \includegraphics[width=0.95\textwidth]{pipeline_architecture.png}
    \caption{End-to-end architecture: an intent is decomposed into sub-tasks,
    which drive two parallel retrieval systems (BraneScript language
    specification and package/dataset context). The retrieved context is
    combined into a prompt that an LLM uses to generate BraneScript, which is
    then syntax-checked, semantically checked, and executed. Failures at either
    check re-enter the prompt-construction stage as feedback for a bounded
    number of retries.}
    \label{fig:pipeline-architecture}
\end{figure}

The pipeline is organised into six conceptual stages:

\begin{enumerate}
    \item \textbf{Intent decomposition.} The raw natural-language intent is
        broken down into a high-level task description: which function(s) need
        to be defined, which package(s) must be imported, which dataset(s) are
        referenced, and what control flow (loops, conditionals) the workflow
        requires.

    \item \textbf{Dual context retrieval.} The decomposed intent feeds two
        independent retrieval systems in parallel: a \emph{language
        specification} retriever, which returns relevant BraneScript syntax
        (function definitions, control-flow constructs, parallel blocks) from a
        dedicated knowledge base; and a \emph{package/dataset} retriever, which
        returns the function signatures and usage context of the specific
        Brane packages and datasets relevant to the intent (e.g., the
        \texttt{healthcare} package and the associated heart-disease dataset).

    \item \textbf{Prompt construction.} The retrieved language specification and
        package/dataset context are combined with the original intent and a
        system prompt into a single generation prompt for the LLM.

    \item \textbf{Generation.} The LLM (a general-purpose or fine-tuned Qwen
        model) produces a candidate BraneScript workflow.

    \item \textbf{Validation.} The candidate is first \emph{syntax-checked}
        against the Brane compiler (a fast, local check that requires no
        execution), and, on success, \emph{semantically checked} against the
        retrieved package context (e.g., that referenced functions and
        arguments actually exist in the imported packages).

    \item \textbf{Execution.} Only workflows that pass both checks are executed
        on the Brane runtime. If either check fails, the error is fed back into
        prompt construction as corrective context, and the pipeline retries
        generation up to a fixed number of attempts before reporting failure to
        the user.
\end{enumerate}

This structure realises the two design principles introduced informally in
Chapter~\ref{s:background} and developed further in
Chapter~\ref{s:design}: \emph{compile-before-execute}, so that no compute
resources are spent on malformed workflows, and \emph{progressive trust}, so
that cheaper, more deterministic mechanisms (the semantic cache, the syntax
checker) are always tried before the more expensive and stochastic ones (LLM
generation, workflow execution).

\section{Components}

At the code level, the architecture in Figure~\ref{fig:pipeline-architecture}
is realised by the following components, each described in detail in
Chapter~\ref{s:implementation}:

\begin{itemize}
    \item \textbf{Pipeline orchestrator} (\texttt{src/pipeline.py}) --- drives
        the six stages above and implements the bounded retry loop.
    \item \textbf{Knowledge base} (\texttt{src/knowledgeBase.py}) --- builds and
        queries the two ChromaDB vector stores that back the dual retrieval
        stage.
    \item \textbf{Semantic cache} (\texttt{src/semantic\_cache.py}) --- a
        ChromaDB-backed cache that short-circuits generation for intents that
        are semantically equivalent to a previously seen one, sitting in front
        of the LLM in the progressive-trust ordering.
    \item \textbf{Brane packages} (\texttt{packages/*}) --- seven
        domain-specific packages (\texttt{healthcare}, \texttt{genomics},
        \texttt{epidemics}, \texttt{statistics}, \texttt{text\_analysis},
        \texttt{data\_masking}, \texttt{datetime}) that expose the typed
        functions the generated BraneScript invokes, and that execute against
        locally registered datasets.
    \item \textbf{Fine-tuning and evaluation infrastructure}
        (\texttt{src/fine\_tuning/train.py}, \texttt{evaluate.py}) --- produces
        and benchmarks the domain-adapted LLMs used as an alternative to the
        base model in the generation stage.
    \item \textbf{Dashboard} (\texttt{frontend/}) --- exposes the pipeline
        interactively, and additionally visualises evaluation results and cache
        statistics.
\end{itemize}

Crucially, the data that packages such as \texttt{healthcare} and
\texttt{genomics} operate on never leaves the site at which it is registered:
the LLM only ever sees the natural-language intent, package function
signatures, and (redacted) dataset metadata during generation; the generated
BraneScript is executed by Brane's runtime against the dataset in place. This
separation between \emph{workflow generation} (which may involve cloud-hosted
or shared LLM infrastructure) and \emph{data access} (which remains local to
the owning site) is the architectural realisation of the privacy and
data-ownership motivation for choosing Brane as the target platform for this
work (Section~\ref{s:background}).

